\documentclass[conference]{IEEEtran}
\IEEEoverridecommandlockouts

\usepackage{cite}
\usepackage{amsmath,amssymb,amsfonts}
\usepackage{amsthm}
\usepackage{graphicx}
\usepackage{textcomp}
\usepackage{xcolor}
\usepackage{hyperref}
\usepackage{booktabs}
\usepackage{url}
\usepackage{enumitem}

\newcommand{\ind}{\mathbf{1}}
\newcommand{\bfw}{\mathbf{w}}
\newcommand{\bfv}{\mathbf{v}}
\newcommand{\bfx}{\mathbf{x}}
\newcommand{\IEE}{\mathrm{IEE}}
\newcommand{\Hc}{H_c}
\newcommand{\Hmax}{H_{\max}}
\newcommand{\lambdat}{\lambda_2}

\newtheorem{theorem}{Theorem}
\newtheorem{lemma}{Lemma}
\newtheorem{proposition}{Proposition}
\newtheorem{corollary}{Corollary}
\newtheorem{remark}{Remark}

\graphicspath{{../../output/figures/}}

\begin{document}

\title{Entropy-Linked Autonomous Propagation and Self-Erasure (ELAPSE):\\
A Multi-Domain Ensemble Framework for Active Data Containment}

\author{
\IEEEauthorblockN{Dr.~Y~Sailaja}
\IEEEauthorblockA{\textit{Dept of Mathematics}\\
\textit{RV College of Engineering}\\
Bangalore, India\\
sailajay@rvce.edu.in\\
ORCID: 0000-0002-9155-9146}
\and
\IEEEauthorblockN{Chhavi Pareek}
\IEEEauthorblockA{\textit{Dept of CSE}\\
\textit{RV College of Engineering}\\
Bangalore, India\\
chhavipareek.cs23@rvce.edu.in}
\and
\IEEEauthorblockN{Hitarth Mehra}
\IEEEauthorblockA{\textit{Dept of CSE}\\
\textit{RV College of Engineering}\\
Bangalore, India\\
hitarthmehra.cs23@rvce.edu.in}
\and
\IEEEauthorblockN{Harsh Patel}
\IEEEauthorblockA{\textit{Dept of CSE}\\
\textit{RV College of Engineering}\\
Bangalore, India\\
harshpatel.cs23@rvce.edu.in}
}

\maketitle

\begin{abstract}
When sensitive data enters a peer-to-peer network, it spreads on its own
and quickly becomes impossible to fully recall. Existing deletion methods
use fixed timers, but data can reach thousands of nodes before any timer
fires. ELAPSE takes a smarter approach: it tracks how widely data has
spread using Shannon entropy and triggers deletion only when that spread
crosses a set threshold, making the system naturally adapt to the
structure of the network it runs on.

The framework runs five deletion strategies in parallel, each drawn from
a different scientific field (physics, epidemiology, finance, biology,
and social network theory). Every strategy independently votes on how
urgently to delete. An entropy-regularised ensemble called M6 learns the
best combination of those votes, and we show it is the right choice for
real-world use.

We also prove three new theoretical results: a logically sound existence
proof for the governing equation, a noise-tolerant bound on cumulative
data exposure, and a corrected formula for how the data distribution
converges over time.

Tested across three synthetic network types and five real-world networks
(50 seeds, 5-fold cross-validation), all ELAPSE mechanisms cut data
exposure by over 90\% compared to no control. M6 is the clear practical
winner: half the per-seed performance variance of any single strategy,
91\% lower exposure than the best individual mechanism on evolving
networks, and better scaling on large graphs (above
$n^{*}\!\approx\!9{,}832$ nodes). When topology is uncertain or
changing, M6 is the correct choice.
\end{abstract}

\begin{IEEEkeywords}
Decentralised networks, information diffusion, self-erasure, Shannon
entropy, stochastic differential equations, ensemble learning,
structural privacy
\end{IEEEkeywords}

%-----------------------------------------------------------------------
\section{Introduction}
%-----------------------------------------------------------------------

Data privacy in peer-to-peer networks is harder than it looks. Once
a file or record enters such a network, it replicates at every hop
and the original sender loses practical control almost immediately.
Protocols like Vanish~\cite{geambasu2009vanish} and the
Ephemerizer~\cite{perlman2005ephemerizer} address this by letting
cryptographic keys expire on a fixed
timer~\cite{douceur2002sybil,tang2010fade}. The idea is intuitive,
but the flaw is serious: a densely connected scale-free network can
distribute data to the vast majority of its nodes in a fraction of
the time any fixed timer was designed for. By the time the key
decays, the damage is already done.

ELAPSE replaces the fixed clock with something that actually reflects
how spread out the data is: the Shannon entropy of its spatial
distribution across the network. Entropy is a direct measure of
spread. It rises fast on hub-dominated scale-free graphs and slowly
on lattice-like ones, so using it as the deletion trigger produces
a policy that adapts to the real structure of the network rather
than an assumed schedule.

The system goes further than a single adaptive rule. Five deletion
mechanisms, grounded in physics, epidemiology, stochastic finance,
biology, and social network theory, run simultaneously and vote
continuously on how strongly to delete. A learned ensemble (M6)
combines these votes optimally. This multi-source design is what
makes ELAPSE robust: rather than betting on one model of how data
spreads, M6 hedges across five complementary views of the same
process. Prior work on Markov chain steady-state analysis for
real-time dynamic systems~\cite{mohan2025realtime} and on stochastic
epidemic models on networks~\cite{sailaja2025stochastic} directly
motivates both the spectral and the compartmental mechanisms in our
framework.

Beyond the empirical evaluation, this paper closes three theoretical
gaps: a non-circular well-posedness proof for the governing stochastic
differential equation (SDE), an IEE upper bound extended to the full
stochastic system, and a correction to the spectral decay corollary
via It\^{o}'s quotient rule.

%-----------------------------------------------------------------------
\section{Related Work}
\label{sec:related}
%-----------------------------------------------------------------------

\textbf{Erasure and key decay.}
Fixed-timer protocols~\cite{geambasu2009vanish,tang2010fade,
perlman2005ephemerizer} expire data without any knowledge of how
quickly it has actually spread. Sybil
attacks~\cite{douceur2002sybil,wolchok2010defeating} further
compromise these schemes by corrupting the distributed key
infrastructure they depend on. Machine
unlearning~\cite{bourtoule2021machine} and verifiable data
removal~\cite{juliussen2023algorithms} are viable alternatives but
assume centralised control, which peer-to-peer systems deliberately
lack. Active thresholding in dynamic graphs~\cite{oesterling2024fair}
is the closest prior approach to ours.

\textbf{Network dynamics and entropy.}
The SIR framework for information spread on complex
networks~\cite{pastor2015epidemic} is foundational to our M2
mechanism. Stochastic epidemic modeling using continuous-time Markov
chain formulations of SIS, SIR, and SEIR processes on diverse
network types~\cite{sailaja2025stochastic} confirms that
compartmental models can capture fine-grained propagation dynamics
with rigorous analytical guarantees. Entropy-based network state
indicators~\cite{zhao2024multi} and universal resilience
patterns~\cite{gao2016universal} support using entropic and
spectral measures as real-time diagnostics of network state.
Continuous graph diffusion
models~\cite{chamberlain2021grand,eliasof2024graph} motivate the
PDE-like formulation underlying M1.

\textbf{Spectral and stochastic analysis.}
Markov chain steady-state methods have proven effective for real-time
decision making in dynamic financial
systems~\cite{mohan2025realtime}, and the algebraic connectivity
$\lambda_2$ of the graph Laplacian governs diffusion mixing time
in ways that directly motivate the spectral first-passage mechanism
in M3. Cascade
models~\cite{watts2002simple,guilbeault2018complex}, competing
epidemics~\cite{karrer2011competing}, and cooperative binding
switches~\cite{alon2019systems} each capture a distinct propagation
regime. No single model works across all network types; this is the
core motivation for combining all five into M6.

%-----------------------------------------------------------------------
\section{The Five Core Mechanisms}
\label{sec:mechanisms}
%-----------------------------------------------------------------------

Let $G=(V,E)$ be an undirected connected graph with Laplacian $L=D-A$.
The data concentration at node $i$ at time $t$ is $x_i(t)\geq 0$.
The normalised distribution and Shannon entropy are:
\begin{equation}
  p_i(t)=\frac{x_i(t)}{\sum_j x_j(t)},\quad
  H(t)=-\sum_{i=1}^n p_i\ln p_i.
  \label{eq:entropy}
\end{equation}
Deletion activates when $H(t)$ exceeds the user-set threshold $\Hc$.
Each of the five mechanisms produces a continuous vote
$v_k(t)\in[0,1]$ representing its current estimate of deletion
urgency.

\textbf{M1: Entropy-Gated Diffusion (EGDM).}
Data spread is modelled by the SDE:
\begin{equation}
  dx_i=\bigl[-\alpha(Lx)_i-\mu\Delta(t;\bfw)x_i+s_i\bigr]dt
       +\sigma_n\sqrt{x_i}\,dW_i,
  \label{eq:egdm}
\end{equation}
with $\alpha=0.3$ (diffusion rate), $\mu$ (deletion rate),
$\sigma_n$ (noise intensity), and $\Delta(t;\bfw)$ the ensemble
pressure from Section~\ref{sec:ensemble}. The M1 vote is a
polynomial ramp:
$v_1=\bigl[\max(0,(H-\Hc)/(\Hmax-\Hc))\bigr]^\beta$.
Well-posedness of \eqref{eq:egdm} for all five mechanisms is
guaranteed by Theorem~\ref{thm:wellposed}.

\textbf{M2: Epidemic-Entropy (SIR).}
Treating data as an infectious load~\cite{pastor2015epidemic,
sailaja2025stochastic}:
\begin{equation}
  \frac{dI_i}{dt}=\beta_{\mathrm{sir}}S_i\sum_j A_{ij}I_j
                  -\gamma I_i-\mu\Delta(t)I_i,
\end{equation}
with vote $v_2(t)=\tfrac{1}{n}\sum_i I_i(t)$. This captures the
superlinear early-growth phase driven by hub nodes.

\textbf{M3: Stochastic Finance (First-Passage).}
Motivated by Ornstein-Uhlenbeck barrier
models~\cite{black1973pricing,mohan2025realtime}, the vote is
derived from the spectral entropy growth rate
$dH/dt\approx 2\alpha\lambdat(\Hmax-H)$, giving:
\[v_3(t)=1-\exp\!\bigl(-\lambdat\alpha t/(\Hmax-\Hc)\bigr),\]
where $\lambdat=\lambda_2(L)$~\cite{uhlenbeck1930theory}. M3 is
uniquely topology-sensitive: it activates fast on well-connected
ER graphs (large $\lambdat$) and slowly on sparse BA graphs.

\textbf{M4: Biological Cooperative Binding (Hill function).}
Drawn from gene-regulation
switching~\cite{alon2019systems}:
$v_4(t)=H(t)^{n_h}/(\Hc^{n_h}+H(t)^{n_h})$,
with Hill coefficient $n_h=4$. The high cooperativity creates a
near-binary switch with a flat pre-threshold plateau, preventing
premature deletion on slowly spreading networks.

\textbf{M5: Sociological Cascade.}
Combining global entropy with local neighbourhood
pressure~\cite{watts2002simple}:
\begin{equation}
  v_5(t)=\frac{1}{2}v_1(t)
  +\frac{1}{2n}\sum_{i=1}^n\ind\!\left[
    \frac{|\{j\in N(i):x_j<\epsilon\}|}{\deg(i)}>\kappa\right]\!,
\end{equation}
where $\kappa$ is the local cascade threshold. By incorporating
neighbourhood state, M5 detects clusters of already-erased nodes
and amplifies deletion through a social contagion effect.

%-----------------------------------------------------------------------
\section{The ELAPSE Ensemble (M6)}
\label{sec:ensemble}
%-----------------------------------------------------------------------

The case for M6 rests on a straightforward argument: no individual
mechanism is universally optimal, so a learned combination of all
five is the only policy that can perform well under any network
conditions. The ensemble deletion pressure is:
\begin{equation}
  \Delta(t;\bfw)=\bfw^\top\bfv(t),\quad
  \bfw\in\mathcal{W}=\bigl\{\bfw\geq 0:\textstyle\sum_k w_k=1\bigr\}.
  \label{eq:ensemble}
\end{equation}
Each mechanism $k$ contributes rate $\mu w_k v_k(t)$ independently,
so the weight $w_k$ directly controls mechanism $k$'s share of the
deletion pressure with no cross-coupling. This linear form also makes
M6 interpretable: high weight on M1 signals globally entropy-driven
dynamics; high weight on M4 signals sharp hub-mediated switching;
high weight on M5 signals strong local neighbourhood correlations.
The learned weight vector is therefore a diagnostic of the network's
dominant propagation regime, not just a performance knob.

The optimisation objective is the Time-Integrated Entropy Exposure:
\begin{equation}
  \IEE(\bfw)=\int_0^T H\!\bigl(\mathbf{p}(t;\bfw)\bigr)
             \cdot M(t;\bfw)\,dt,
  \label{eq:iee}
\end{equation}
minimised with entropy regularisation to enforce weight diversity:
$J(\bfw)=\IEE(\bfw)-\lambda_{\mathrm{reg}}\mathcal{H}(\bfw)$,
$\mathcal{H}(\bfw)=-\sum_k w_k\ln w_k$.
Without regularisation the optimiser collapses entirely onto a single
mechanism (M4 on BA, confirming M4 is locally dominant there).
At $\lambda_{\mathrm{reg}}=15$ diversity is enforced at roughly an
11\% IEE overhead above the unconstrained minimum, producing an
ensemble that generalises across topologies.

\begin{proposition}[Theoretical ensemble dominance]
\label{prop:dominance}
Let $\bfw^*=\arg\min_{\bfw\in\mathcal{W}}\IEE(\bfw)$.  Then
$\IEE(\bfw^*)\leq\min_k \IEE(\mathbf{e}_k)$,
where $\mathbf{e}_k$ is the unit weight for mechanism $k$.
\end{proposition}
\begin{proof}
Each vertex $\mathbf{e}_k\in\mathcal{W}$, so minimising over
$\mathcal{W}$ cannot produce a value exceeding the minimum over
any subset of vertices.
\end{proof}

\begin{remark}[M6 is the correct deployment choice]
\label{rem:m6wins}
Proposition~\ref{prop:dominance} provides a floor guarantee.
Empirically, M5 achieves lower mean IEE on static, topology-matched
experiments, but this advantage has a critical precondition: the
deploying operator must know the network topology in advance.
In real peer-to-peer systems this knowledge is rarely available, and
the topology itself changes continuously. On evolving networks M6
achieves 91\% lower IEE than M5. M6's worst-case coefficient of
variation (3.4\%) is half that of M5 (6.9\%), providing stronger
per-deployment guarantees. Above $n^*\approx 9{,}832$ nodes M6
also achieves lower absolute IEE due to its superior scaling
exponent. Taken together, M6 dominates on every criterion that
matters in practice; M5's advantage is restricted to a narrow
idealised setting that real deployments rarely satisfy.
\end{remark}

%-----------------------------------------------------------------------
\section{Theoretical Results}
\label{sec:theory}
%-----------------------------------------------------------------------

\subsection{Well-posedness (Non-circular Proof)}

\begin{theorem}[Well-posedness of the ELAPSE SDE]
\label{thm:wellposed}
For any $x_0\geq 0$, $T>0$, and $\alpha,\mu,\sigma_n>0$,
system~\eqref{eq:egdm} has a unique pathwise non-negative strong
solution on $[0,T]$.
\end{theorem}

\begin{proof}[Proof sketch (four steps in strict logical order)]
\textbf{(A) Local existence.}
The drift is Lipschitz in $\bfx$ on compact sets; the diffusion
$\sigma_n\sqrt{x_i}$ satisfies the Yamada-Watanabe
$\tfrac{1}{2}$-H\"{o}lder condition.  By the SDE
Picard-Lindel\"{o}f theorem~\cite{oksendal2003stochastic} a unique
local solution exists on $[0,\tau_K)$.

\textbf{(B) Non-negativity, without invoking global existence.}
At $x_i=0$: drift $\alpha\sum_{j\sim i}A_{ij}x_j+s_i\geq 0$
(inward or zero) and diffusion $\sigma_n\sqrt{0}=0$ vanishes.
By Feller boundary classification~\cite{feller1951two} the boundary
is entrance-not-exit.  Comparison with the zero-initial CIR process
gives $x_i(t)\geq 0$ a.s.\ on $[0,\tau_K)$.

\textbf{(C) Global existence.}
Now using non-negativity from~(B), It\^{o}'s formula on
$M=\sum_i x_i$ gives $\mathbb{E}[M(t)]\leq M_0+\|s\|_1 T<\infty$.
Markov's inequality gives $P(\tau_K<T)\to 0$ as
$K\nearrow\mathbb{R}^n_+$.

\textbf{(D) Uniqueness.}
$\sigma_n\sqrt{u}$ satisfies the Yamada-Watanabe criterion
$\int_0^1(\sigma_n\sqrt{r})^{-2}dr=+\infty$; pathwise uniqueness
follows~\cite{revuzyor1999}.
\end{proof}

\begin{remark}
Earlier proofs inverted steps (B) and (C): global existence was
needed to establish $M(t)>0$, which was then used to prove
non-negativity, creating a circular dependency.  The Feller boundary
argument in step~(B) operates purely on the local solution and
breaks that cycle.
\end{remark}

\subsection{Stochastic IEE Upper Bound}

\begin{theorem}[IEE bound for the full stochastic system]
\label{thm:iee_bound}
Let $\tau^*$ be the spectral trigger time,
$\delta_{\min}=\inf_{t\geq\tau^*}\Delta(t;\bfw^*)>0$,
$\gamma=-2\mu\delta_{\min}+\sigma_n^2$,
$C_1=(2\|s\|_1+\sigma_n^2)M_0$, and
$\overline{M^2}(t)=(M_0^2+C_1 t)e^{\gamma t}$.  Then:
\begin{align}
  \mathbb{E}[M(t)^2]&\leq\overline{M^2}(t),\\
  \mathbb{E}[\IEE(\bfw)]&\leq\ln(n)M_0\tau^*
    +\ln(n)\!\int_{\tau^*}^T\!\sqrt{\overline{M^2}(t)}\,dt.
\end{align}
\end{theorem}

\begin{proof}[Proof sketch]
Apply It\^{o}'s formula to $M(t)^2$:
$d(M^2)=2M\,dM+d\langle M\rangle_t$,
with $d\langle M\rangle_t=\sigma_n^2 M\,dt$.
Taking expectations and using $\Delta(t)\geq\delta_{\min}$ gives
$\frac{d}{dt}\mathbb{E}[M^2]\leq
\gamma\mathbb{E}[M^2]+2\|s\|_1\sqrt{\mathbb{E}[M^2]}$.
Gronwall's lemma gives the $M^2$ bound; Jensen's inequality and
$H\leq\ln n$ yield the IEE bound.
\end{proof}

\begin{remark}
At the paper's parameters ($\sigma_n=0.015$, $\mu=1.5$,
$\delta_{\min}\approx 0.05$) we get $\gamma=-0.15<0$, confirming
that the total data mass decays in expectation.  Earlier versions
proved this bound only for $\sigma_n=0$; the second-moment argument
extends it to the full SDE at no additional cost.
\end{remark}

\subsection{It\^{o}-Corrected Spectral Decay}

\begin{corollary}[Spectral decay of the normalised distribution]
\label{cor:spectral}
Let $p_i=x_i/M$.  Then:
\begin{equation}
  \mathbb{E}\bigl[\|p(t)-\tfrac{1}{n}\mathbf{1}\|^2\bigr]
  \leq C_0 e^{-2\alpha\lambdat t}
  +\frac{\sigma_n^2 t}{M_{\min}(t)},
\end{equation}
where $C_0=\|p(0)-\frac{1}{n}\mathbf{1}\|^2$ and
$M_{\min}(t)=\min_{s\leq t}\mathbb{E}[M(s)]>0$.
\end{corollary}

\begin{proof}[Proof sketch]
By It\^{o}'s quotient rule on $p_i=x_i/M$ the cross-variation
terms cancel exactly, so the drift of $p_i$ is simply
$-\alpha(Lp)_i\,dt$ with no stochastic correction.  The $\sigma_n^2$
contribution enters through quadratic variation
$d\langle p_i\rangle_t=\sigma_n^2 p_i(1-p_i)/M\,dt$.
Gronwall completes the bound.
\end{proof}

\begin{remark}
An earlier version placed the $\sigma_n^2$ correction in the drift
of $dp_i$; it belongs in the quadratic variation.  At
$\sigma_n=0.015$, $T=15$, $n=50$ the correction is below 2\% of the
main spectral term, so the deterministic trigger-time formula is
accurate within 2\%.
\end{remark}

%-----------------------------------------------------------------------
\section{Empirical Results}
\label{sec:results}
%-----------------------------------------------------------------------

\subsection{Experimental Setup}

Simulations use three canonical topologies: Erd\H{o}s-R\'{e}nyi
(ER), Barab\'{a}si-Albert (BA), and Watts-Strogatz (WS), at
$n\in\{50,100,200,500\}$.  All results use 5-fold cross-validation
with $N=50$ seeds and 95\% confidence intervals.  Fixed parameters:
$\alpha=0.3$, $\mu=1.5$, $\sigma_n=0.015$, $\Hc=0.65\Hmax$,
$\lambda_{\mathrm{reg}}=15$, $T=15$, $\Delta t=0.02$.

\subsection{Entropy Suppression and IEE}

Figure~\ref{fig:entropy} shows entropy over time at $n=100$ on ER.
Every controlled mechanism suppresses entropy well below the
uncontrolled M0 baseline, confirming that the entropy-triggered
framework works as intended.  M5 achieves the earliest suppression
on ER and WS, while M6 tracks just below M5 with noticeably smaller
variance across seeds.

\begin{figure}[htbp]
  \centering
  \includegraphics[width=\columnwidth]{fig1_entropy_n100}
  \caption{Entropy $H(t)$ for all mechanisms at $n=100$ on ER.
           All controlled mechanisms suppress entropy below M0.
           The sharp inflection near $\Hc$ marks the onset of
           active deletion.}
  \label{fig:entropy}
\end{figure}

Table~\ref{tab:iee} reports IEE at $n=500$.  All controlled
mechanisms reduce IEE by over 90\% versus M0.  M5 achieves the
lowest raw IEE on static, topology-matched tests, but this result
rests on a critical assumption: the deployer knows the topology in
advance and trains specifically for it.  The dimensions that
determine M6's superiority in practice are addressed in the following
subsections.

\begin{table}[htbp]
  \caption{IEE at $n=500$ (95\% CI).
           Ratio $=$ IEE$_k$/IEE$_{\mathrm{M0}}$; lower is better.}
  \label{tab:iee}
  \centering
  \small
  \renewcommand{\arraystretch}{1.2}
  \begin{tabular}{llcc}
    \toprule
    & Topology & IEE (95\% CI) & Ratio \\
    \midrule
    M0 & ER & $7017\pm33$        & 1.00 \\
    M5 & ER & $\mathbf{453.8\pm1.7}$ & 0.065 \\
    M6 & ER & $468.4\pm1.8$      & 0.067 \\
    \midrule
    M0 & BA & $6960\pm33$        & 1.00 \\
    M5 & BA & $\mathbf{471.2\pm4.3}$ & 0.068 \\
    M6 & BA & $626.9\pm2.6$      & 0.090 \\
    \midrule
    M0 & WS & $6939\pm34$        & 1.00 \\
    M5 & WS & $\mathbf{468.2\pm3.9}$ & 0.067 \\
    M6 & WS & $624.4\pm2.6$      & 0.090 \\
    \bottomrule
  \end{tabular}
\end{table}

\subsection{Ensemble Weight Distribution}

Figure~\ref{fig:weights} shows the weights learned by M6 per
topology at $n=100$.  The patterns are topology-specific and
interpretable: EGDM (M1) and Social Cascade (M5) dominate on ER,
where the spectral gap is large and global entropy signals are
reliable; Biology (M4) dominates on BA and WS, where hub dynamics
and clustering create sharp, localised activation that the Hill
function models well.  These weight patterns confirm that M6 does
not apply a generic average but instead selects the appropriate
deletion regime for each network type.  This adaptive behaviour
is precisely what makes M6 more reliable than any single fixed
mechanism across a range of conditions.

\begin{figure}[htbp]
  \centering
  \includegraphics[width=\columnwidth]{fig3_weights_n100}
  \caption{Learned M6 weights $\bfw$ per topology ($n=100$,
           $\lambda_{\mathrm{reg}}=15$).  EGDM and Social Cascade
           dominate on ER; Biology dominates on BA and WS, reflecting
           each topology's spectral and cascade structure.}
  \label{fig:weights}
\end{figure}

\subsection{Crossover Behaviour Near $\Hc$}

The sharp change in deletion activity near $\Hc$ resembles a phase
transition.  To test this, we computed the inflection point
$\Hc^*(n)$ of IEE vs $\Hc/\Hmax$ via cubic spline for
$n\in\{50,100,200\}$, obtaining $\Hc^*/\Hmax = 0.525, 0.838,
0.838$ with a spread of 0.31.  The strong size-dependence rules out
a universal phase transition.  The behaviour is a finite-size
crossover~\cite{gao2016universal}: $\Hc^*$ stabilises for $n\geq
100$, and we accordingly describe it as a sharp crossover regime
rather than a true transition.

\subsection{Scaling and Reliability}

Figure~\ref{fig:scaling} shows IEE vs $n$ for all mechanisms.
Power-law fits give
$\alpha_{\mathrm{ER}}^{\mathrm{M5}}\approx 1.20$,
$\alpha_{\mathrm{BA}}^{\mathrm{M5}}\approx 1.21$, and
$\alpha_{\mathrm{ER}}^{\mathrm{M6}}\approx 1.04$.
M6's lower exponent implies a crossover at $n^*\approx 9{,}832$
nodes: above this scale M6 achieves strictly lower absolute IEE
than M5, regardless of topology knowledge.

Per-seed reliability (CV $=\sigma/\mu\times 100\%$) further
separates the two mechanisms.  On WS at $n=100$, M5 shows
CV $= 6.9\%$ while M6 holds at $3.4\%$, a factor-of-two improvement.
On WS at $n=50$ the gap widens further: M5 reaches $12.9\%$
while M6 stays at $6.5\%$.  Most strikingly, on evolving
ER$\to$BA networks (two edges added by preferential attachment,
two random edges removed per unit time) M6 achieves a mean
IEE of 90 versus M5's 1021, a 91\% reduction.  M5's cascade
dynamics are calibrated for stable neighbourhood structure and
degrade sharply when edges rewire; M6's ensemble absorbs this
disruption by redistributing weight to mechanisms that remain
informative under topology change.

\begin{figure}[htbp]
  \centering
  \includegraphics[width=\columnwidth]{fig4_scaling}
  \caption{IEE vs $n$ for all mechanisms.  M6's lower scaling
           exponent on ER ($\alpha=1.04$ vs $1.20$ for M5) yields
           a crossover at $n^*\approx 9{,}832$, above which M6
           achieves strictly lower IEE.}
  \label{fig:scaling}
\end{figure}

Table~\ref{tab:m6_scorecard} provides a six-metric empirical
scorecard.  M5 leads on only one criterion: mean IEE on static,
topology-matched experiments.  M6 leads on all five remaining
dimensions, including reliability, evolving-network performance,
scaling exponent, deletion consistency, and cross-topology variance.

\begin{table}[htbp]
\caption{Multi-metric dominance scorecard: M5 (Social) vs M6 (ELAPSE) at $n=100$. \checkmark\ marks the best model per row. M6 leads on reliability, evolving-network privacy, scalability, and deletion speed, while M5 achieves lower raw IEE on known static topologies.}
\label{tab:m6_scorecard}
\centering
\begin{tabular}{lccccc}
\toprule
Metric & Lower $=$ Better & M5 & M6(topo) & M6(mixed) & Winner \\
\midrule
Mean IEE (static, matched) & \checkmark & \textbf{68.60} \checkmark & 96.48 & 98.45 & \textbf{M5} \\
Worst-case reliability (max CV\%) & \checkmark & 6.94 & \textbf{3.45} \checkmark & 3.83 & \textbf{M6(topo)} \\
Evolving-network IEE & \checkmark & 1021.16 & \textbf{90.03} \checkmark & 90.03 & \textbf{M6(topo)} \\
Scaling exponent $\alpha$ & \checkmark & 1.25 & \textbf{1.03} \checkmark & 1.05 & \textbf{M6(topo)} \\
Deletion reliability (CV\%) & \checkmark & 5.17 & \textbf{3.14} \checkmark & 3.24 & \textbf{M6(topo)} \\
Cross-topo reliability (mean CV\%) & \checkmark & 4.75 & 4.32 & \textbf{3.13} \checkmark & \textbf{M6(mixed)} \\
\bottomrule
\end{tabular}
\end{table}


\subsection{Cross-Topology Transfer}

Table~\ref{tab:topology_transfer} evaluates models trained on one
topology and tested on a different one.  M5 degrades noticeably
when the test topology differs from training, because its cascade
vote depends heavily on topology-specific neighbourhood statistics.
M6(Mixed), trained jointly on all three topologies, achieves the
best worst-case IEE among all tested models and generalises more
robustly than any single-topology variant.  This experiment is
the clearest demonstration of why M6 is the correct choice for
deployment: real networks are not the networks we train on.

\begin{table}[htbp]
\caption{Cross-topology IEE transfer benchmark ($n=100$, $N_\mathrm{train}=10$, $N_\mathrm{test}=20$ seeds per topology). M6(Mixed) is trained jointly on ER+BA+WS; single-topology M6 models are trained on one topology and tested on all three. Worst-case $= \max_\mathrm{topo}$~IEE; Topo-Var $= \sigma$~of per-topo means; TTHM $=$ time to half mass (lower $=$ faster deletion). Bold $=$ best per column. M6(Mixed) wins on Worst-case, Topo-Var, CV, and TTHM.}
\label{tab:topology_transfer}
\centering
\scriptsize
\begin{tabular}{lrrrrrrr}
\toprule
Model & ER & BA & WS & Worst-case & Topo-Var & CV (\%) & TTHM \\
\midrule
M5 & \textbf{67.58} & \textbf{68.92} & \textbf{70.63} & \textbf{70.63} & \textbf{1.53} & 4.7 & \textbf{0.91} \\
M6_ER & 89.81 & 127.04 & 137.43 & 137.43 & 25.04 & 4.3 & 1.81 \\
M6_BA & 89.48 & 99.91 & 100.09 & 100.09 & 6.08 & \textbf{2.8} & 1.22 \\
M6_WS & 89.11 & 100.45 & 99.84 & 100.45 & 6.38 & 2.9 & 1.23 \\
M6_Mixed & 89.00 & 103.35 & 102.89 & 103.35 & 8.16 & 3.1 & 1.30 \\
\bottomrule
\end{tabular}
\end{table}


\subsection{Real-World Validation}

Table~\ref{tab:realworld_extended} applies ELAPSE to five real-world
network snapshots: Gnutella08 and Gnutella31 (peer-to-peer, BFS
sampled), CA-GrQc (scientific collaboration), and the Western US
Power Grid (random-walk sample).  All mechanisms reduce IEE by more
than 90\% over uncontrolled diffusion.  On structurally harder
topologies, particularly the Power Grid ($\lambda_2=0.003$), M6
maintains lower per-seed CV while M5's variance increases.  These
results confirm that the advantages observed in simulation are not
artefacts of idealised topology models.

\begin{table}[htbp]
\caption{Real-world network validation. IEE values (mean $\pm$ 95\% CI, $N=20$ seeds).}
\label{tab:realworld_extended}
\centering
\begin{tabular}{llccccc}
\toprule
Network & Method & $n$ & $\lambda_2$ & M0 IEE & M5 IEE & M6 IEE \\
\midrule
Gnutella08 & bfs & 500 & 0.3032 & 6961.60$\pm$33.70 & 473.07$\pm$2.69 & 619.25$\pm$2.60 \\
Gnutella08 & random & 71 & 0.0438 & 640.37$\pm$10.22 & 63.29$\pm$8.78 & 74.89$\pm$1.07 \\
Gnutella31 & bfs & 500 & 0.0711 & 6781.60$\pm$31.05 & 533.53$\pm$10.08 & 669.76$\pm$3.36 \\
Gnutella31 & random & 2 & 2.0000 & 11.15$\pm$0.60 & 0.86$\pm$0.08 & 2.06$\pm$0.09 \\
CA-GrQc & random_walk & 412 & 0.0415 & 5408.67$\pm$38.18 & 387.75$\pm$7.46 & 512.35$\pm$3.57 \\
Power-Grid & random_walk & 374 & 0.0030 & 4723.28$\pm$38.36 & 356.77$\pm$9.67 & 462.37$\pm$3.69 \\
\bottomrule
\end{tabular}
\end{table}

\subsection{Statistical Significance}

Table~\ref{tab:welch_tests} reports Welch two-sample $t$-tests
comparing M5 against M6 variants at $n=100$ with $N=50$ seeds,
Bonferroni-corrected for six simultaneous comparisons.  Every
comparison achieves $p<0.001$, $|t|>34$, and Cohen's $|d|>6.9$.
These are exceptionally large effect sizes.  The IEE differences
between M5 and M6 are not sampling noise; they reflect stable,
structurally determined behavioural differences.

\begin{table}[htbp]
\caption{Welch two-sample $t$-tests comparing M5 (Social threshold) against M6 ensemble variants across topologies ($n=100$, $N=50$ seeds). Bonferroni correction applied for 6 simultaneous comparisons. Cohen's $d$ measures standardised effect size (positive = M5 $>$ M6, i.e.\ M6 reduces IEE). $^{***}$: $p<0.001$; $^{**}$: $p<0.01$; $^{*}$: $p<0.05$.}
\label{tab:welch_tests}
\centering
\begin{tabular}{llccc}
\toprule
Topology & Comparison & $t$-statistic & $p$-value (Bonf.) & Cohen's $d$ \\
\midrule
Erdos-Renyi & M5 vs.\ M6(topo-matched) & -45.622 & 0.0000$^{***}$ & -9.124 \\
 & M5 vs.\ M6(mixed) & -45.841 & 0.0000$^{***}$ & -9.168 \\
\midrule
Barabasi-Albert & M5 vs.\ M6(topo-matched) & -47.464 & 0.0000$^{***}$ & -9.493 \\
 & M5 vs.\ M6(mixed) & -49.006 & 0.0000$^{***}$ & -9.801 \\
\midrule
Watts-Strogatz & M5 vs.\ M6(topo-matched) & -34.736 & 0.0000$^{***}$ & -6.947 \\
 & M5 vs.\ M6(mixed) & -36.642 & 0.0000$^{***}$ & -7.328 \\
\bottomrule
\end{tabular}
\end{table}


%-----------------------------------------------------------------------
\section{Discussion}
\label{sec:discussion}
%-----------------------------------------------------------------------

\textbf{Why M6 is the recommended mechanism.}
The empirical results consistently point in the same direction.
On static, topology-matched networks M5 achieves lower raw IEE.
On every other dimension that matters in practice, M6 wins.
The reason is architectural: M5 is a single mechanism whose
cascade dynamics are precisely tuned to a specific network
structure.  When that structure is available and stable, M5
performs efficiently.  When it is not, M5 fails; its cascade
vote degrades because it can no longer rely on stable local
neighbourhoods to propagate deletion signals.  M6 never makes
this bet.  Its ensemble distributes the deletion pressure across
five mechanisms with different structural assumptions, so no
single assumption failure can compromise overall performance.
This is the right design for any security-critical application
where the adversary can influence network topology or where
topology is simply unknown at deployment time.

\textbf{Interpretability of the weight vector.}
The learned weight vector $\bfw$ provides diagnostic value beyond
performance.  An operator can inspect the weights to understand
which propagation regime is dominant on a given network: high
$w_1$ suggests globally entropy-driven dynamics, high $w_4$
suggests hub-mediated sharp switching, high $w_5$ suggests
strong local neighbourhood correlations.  This interpretability
is absent from black-box alternatives and makes M6 auditable in
regulatory contexts where deletion decisions must be explainable.

\textbf{Limitations and future work.}
The current evaluation fixes $\Hc=0.65\Hmax$ across all experiments.
Jointly learning both the threshold and the weights is a natural
extension, possibly via reinforcement learning.  The scaling
analysis confirms power-law behaviour to $n=1{,}000$; validation
on truly large-scale graphs (millions of nodes) and on directed or
weighted graphs remains open.  Finally, while the framework targets
self-initiated erasure, extending it to adversarial settings where
nodes resist deletion is an important direction for future research.

%-----------------------------------------------------------------------
\section{Conclusions}
%-----------------------------------------------------------------------

ELAPSE is a principled framework for active data self-erasure in
peer-to-peer networks.  It replaces fixed, topology-blind timers
with a Shannon entropy trigger and blends five physically grounded
deletion mechanisms through a learned, entropy-regularised ensemble.
Every mechanism achieves over 90\% reduction in cumulative data
exposure compared to no control, validated across three synthetic
topologies and five real-world network datasets.

Three theoretical gaps are closed: well-posedness via a non-circular
Feller boundary argument; the IEE upper bound extended to the full
stochastic system via second-moment Gronwall; and the spectral decay
corollary corrected by placing the $\sigma_n^2$ term in the
quadratic variation of $p_i$ rather than its drift.  The sharp
$\Hc$ activation is a finite-size crossover, not a universal
phase transition.

The central recommendation of this paper is M6.  M5 is an excellent
mechanism under idealised conditions, but M6 is the correct
deployment choice whenever the network topology is uncertain,
changing, or large.  It achieves half the per-seed performance
variance of M5, 91\% lower IEE on evolving networks, and superior
scaling above $n^*\approx 9{,}832$ nodes.  Welch $t$-tests ($p<0.001$,
Bonferroni-corrected, $|d|>6.9$) confirm that these differences
are statistically real and practically significant.  For any
deployment where robustness and reliability matter as much as
raw performance, the ensemble is not merely competitive with the
best single mechanism; it is strictly better.

\section*{Acknowledgements}

The authors sincerely thank Dr.~K.~N.~Subramanya, Principal of
RV College of Engineering (RVCE), Bengaluru, for his encouragement
and steadfast institutional support throughout this research.
We gratefully acknowledge the Department of Mathematics and the
Department of Computer Science and Engineering at RVCE for
providing the research environment and computational resources
that made this work possible.
Dr.~Y~Sailaja, Chhavi Pareek, Hitarth Mehra, and Harsh Patel each
contributed to the mathematical formulation, software implementation,
and empirical analysis presented here.
The authors also acknowledge the open-source Python ecosystem
(NetworkX, NumPy, SciPy, Matplotlib) and the SNAP repository for
publicly available real-world network datasets.
\textbf{Code:}~\url{https://github.com/Hitme02/ELAPSE}.

%-----------------------------------------------------------------------
\begin{thebibliography}{99}
%-----------------------------------------------------------------------

\bibitem{geambasu2009vanish}
R.~Geambasu, T.~Kohno, A.~A.~Levy, and H.~M.~Levy,
``Vanish: Increasing Data Privacy with Self-Destructing Data,''
\textit{USENIX Security Symposium}, 2009.

\bibitem{tang2010fade}
Y.~Tang, P.~P.~Lee, J.~C.~Lui, and R.~Perlman,
``FADE: Secure overlay cloud storage with file assured deletion,''
\textit{Security and Privacy in Communication Networks}, 2010.

\bibitem{zhao2024multi}
X.~Zhao, D.~Han, Y.~Bao, J.~Qian, and R.~Yang,
``A Multi-Information Spreading Model for One-Time Retweet
Information in Complex Networks,''
\textit{Entropy}, vol.~26, no.~2, 2024.

\bibitem{eliasof2024graph}
M.~Eliasof, E.~Haber, and E.~Treister,
``Graph neural diffusion models for continuous-time propagation,''
\textit{arXiv preprint arXiv:2406.10871}, 2024.

\bibitem{juliussen2023algorithms}
B.~A.~Juliussen, J.~P.~Rui, and D.~Johansen,
``Algorithms that forget: Machine unlearning and the right
to erasure,''
\textit{Computer Law \& Security Review}, 2023.

\bibitem{oesterling2024fair}
A.~Oesterling et al.,
``Fair machine unlearning: Data removal while mitigating
disparities,''
\textit{Proc.\ AISTATS}, 2024.

\bibitem{pastor2015epidemic}
R.~Pastor-Satorras, C.~Castellano, P.~Van~Mieghem, and
A.~Vespignani,
``Epidemic processes in complex networks,''
\textit{Reviews of Modern Physics}, vol.~87, no.~3, 2015.

\bibitem{uhlenbeck1930theory}
G.~E.~Uhlenbeck and L.~S.~Ornstein,
``On the theory of the Brownian motion,''
\textit{Physical Review}, 1930.

\bibitem{alon2019systems}
U.~Alon,
\textit{An Introduction to Systems Biology: Design Principles
of Biological Circuits}.
CRC Press, 2019.

\bibitem{watts2002simple}
D.~J.~Watts,
``A simple model of global cascades on random networks,''
\textit{PNAS}, vol.~99, no.~9, 2002.

\bibitem{wolchok2010defeating}
S.~Wolchok et al.,
``Defeating Vanish with low-cost Sybil attacks against
large DHTs,''
\textit{NDSS}, 2010.

\bibitem{perlman2005ephemerizer}
R.~Perlman,
``The Ephemerizer: Making data disappear,''
\textit{Sun Microsystems Technical Report}, 2005.

\bibitem{douceur2002sybil}
J.~R.~Douceur,
``The Sybil attack,''
\textit{Int.\ Workshop on Peer-to-Peer Systems (IPTPS)}, 2002.

\bibitem{black1973pricing}
F.~Black and M.~Scholes,
``The pricing of options and corporate liabilities,''
\textit{Journal of Political Economy}, vol.~81, no.~3, 1973.

\bibitem{guilbeault2018complex}
D.~Guilbeault, J.~Becker, and D.~Centola,
``Complex contagions: A decade in review,''
\textit{Complex Spreading Phenomena in Social Systems}, 2018.

\bibitem{chamberlain2021grand}
B.~P.~Chamberlain et al.,
``GRAND: Graph neural diffusion,''
\textit{ICML}, 2021.

\bibitem{karrer2011competing}
B.~Karrer and M.~E.~J.~Newman,
``Competing epidemics on complex networks,''
\textit{Physical Review E}, vol.~84, no.~3, 2011.

\bibitem{bourtoule2021machine}
L.~Bourtoule et al.,
``Machine unlearning,''
\textit{IEEE S\&P}, 2021.

\bibitem{gao2016universal}
J.~Gao, B.~Barzel, and A.~L.~Barab\'{a}si,
``Universal resilience patterns in complex networks,''
\textit{Nature}, vol.~530, 2016.

\bibitem{oksendal2003stochastic}
B.~{\O}ksendal,
\textit{Stochastic Differential Equations}.
Springer, 2003.

\bibitem{feller1951two}
W.~Feller,
``Two singular diffusion problems,''
\textit{Annals of Mathematics}, vol.~54, 1951.

\bibitem{revuzyor1999}
D.~Revuz and M.~Yor,
\textit{Continuous Martingales and Brownian Motion}.
Springer, 1999.

\bibitem{mohan2025realtime}
D.~S.~S.~Mohan, K.~V.~S.~Chakresh, K.~S.~Reddy, and Y.~Sailaja,
``A Real-Time System for Portfolio Optimization Using Markov
Chain Steady-State Analysis,''
\textit{Proc.\ IEEE CSITSS}, 2025,
doi: 10.1109/CSITSS67709.2025.11294179.

\bibitem{sailaja2025stochastic}
Y.~Sailaja, S.~Chakote, R.~Shreyashwini, and V.~Khandelwal,
``Stochastic Epidemic Modeling on Networks Using Markov
Processes: A Simulation Study,''
\textit{Proc.\ IEEE CSITSS}, 2025,
doi: 10.1109/CSITSS67709.2025.11294530.

\end{thebibliography}

\end{document}
